\chapter[Método]{Método}
\label{cap:metodo}

O trabalho será conduzido como uma pesquisa aplicada, quantitativa e experimental em Visão Computacional. O método consiste na implementação e na avaliação controlada de uma pipeline para detecção, rastreamento multiobjeto, estimação de movimento aparente e predição de trajetórias de espermatozoides em vídeos microscópicos. A hipótese será examinada por experimentos reproduzíveis que comparem modelos equivalentes com e sem características locais de fluxo óptico.

\section{Produtos e dados da pesquisa}

Os produtos derivados serão uma coleção de dados organizada com divisões experimentais documentadas; uma pipeline modular em Python; modelos treinados; arquivos contendo caixas, identificadores, trajetórias e campos vetoriais; e uma análise comparativa apresentada por métricas, tabelas e visualizações. O código e os modelos serão artefatos necessários à realização dos experimentos, enquanto o produto científico principal será a evidência sobre o efeito do fluxo óptico na predição de trajetórias.

Os vídeos microscópicos utilizados serão exclusivamente os do VISEM, que contém 85 vídeos, um por participante \cite{haugen2019visem}. O VISEM-Tracking fornece caixas, classes e identificadores persistentes por quadro para trechos iniciais de 30 segundos de 20 desses vídeos \cite{thambawita2023visemtracking}. Os trechos anotados dessas 20 unidades constituirão o gabarito para desenvolvimento e avaliação quantitativa. Os outros 65 vídeos serão processados somente após o congelamento das configurações, como aplicação e análise qualitativa, pois não possuem tracking manual que permita declarar acurácia. Os 502 clipes sem rótulos extraídos dos mesmos vídeos serão tratados como derivados, e não como amostras independentes. Intervalos sem arquivo de anotação no vídeo 23 serão excluídos de treino e métricas, nunca interpretados como quadros negativos. A avaliação não pressupõe contracorrente induzida nem inclui uma segunda coleção experimental com tubo.

\section{Levantamento, geração e organização dos dados}

Para cada vídeo serão registrados resolução, taxa de quadros, duração, escala espacial quando disponível e condições de aquisição. Os quadros serão extraídos preservando a ordem temporal e vinculados a um identificador único. O pré-processamento será limitado a operações cuja influência possa ser controlada, como conversão de cor, normalização de intensidade, correção moderada de contraste e redução de ruído. As anotações serão convertidas para uma representação comum, e as coordenadas das trajetórias serão obtidas a partir do centro das caixas.

Será adotada uma separação fixa por vídeo em 12 unidades de treinamento, quatro de validação e quatro de teste. O treinamento servirá para estimar parâmetros; a validação, para escolher hiperparâmetros e critérios de parada; e a bateria confirmatória no teste ocorrerá somente após o congelamento da configuração. A confirmação em cinco folds dependerá de um desenho explícito que separe ajuste e seleção dos vídeos avaliados em cada fold. O piloto anterior examinou frames dos quatro vídeos hoje reservados ao teste; portanto, esse holdout não é completamente cego. Repetir uma configuração escolhida nos cinco grupos ou redistribuir os vídeos não restaura a independência perdida por essa exposição histórica, que será declarada como limitação. Arquivos de configuração registrarão divisões, sementes 42, 123 e 2026, versões das bibliotecas, hardware e parâmetros resolvidos.

\section{Detecção e rastreamento multiobjeto}

A detecção começará com baselines clássicos separados: threshold fixo com morfologia, Otsu, threshold adaptativo, Blob Detector, MOG2, KNN e Watershed. Uma variante híbrida com correção de fundo e normalização local será avaliada sem substituir os clássicos puros. Por fim, será treinado um detector da família YOLO com as anotações do VISEM-Tracking. A métrica principal será o F1 de localização dos indivíduos anotados nas classes 0 e 2, obtido por matching Húngaro um-para-um entre centros, com tolerância de 10 pixels na resolução original. Todas as previsões participam da associação, independentemente da classe prevista. Após a associação, previsões residuais a até esse raio de qualquer indivíduo anotado permanecem falsos positivos; das restantes, são ignoradas somente aquelas cujo centro está dentro de uma caixa de agrupamento da classe 1. Os indivíduos anotados dentro dessas caixas continuam sendo avaliados. A regra será recalculada integralmente nas sensibilidades de 15 e 20 pixels. Uma avaliação complementar considerará todos os objetos anotados, incluindo cada agrupamento como um objeto, sem interpretá-lo como uma única célula. Serão informadas as contagens brutas, avaliadas e ignoradas, a cobertura espacial das regiões de agrupamento, precisão, revocação, erro de centro, erro de contagem, latência e memória. A MAE de contagem utilizará o mesmo universo de quadros anotados para todos os detectores. As caixas manuais de agrupamento serão usadas exclusivamente na avaliação, sem filtrar entradas ou saídas do detector e do rastreador. As coordenadas, dimensões e valores de confiança serão exportados com a precisão de ponto flutuante usada no processamento, permitindo reconstruir as associações e os erros a partir dos arquivos de saída; arredondamentos serão restritos à apresentação. A tolerância é uma convenção operacional de localização em relação ao centro geométrico da caixa anotada, não uma estimativa de incerteza dos anotadores ou do centro anatômico exato. mAP@0,5, mAP@0,5:0,95 e a análise nas três classes serão secundários e específicos do YOLO. O método de Zhang et al. \cite{zhang2024spermyolo} será utilizado como referência do domínio, considerando as diferenças de dados e ambiente computacional.

Para a triagem prospectiva do threshold fixo, foi registrado um espaço de 171 configurações: limiares de 0 a 255 em passos de 16, incluindo adicionalmente 190 e 200 como referências históricas, com abertura e fechamento de zero, uma ou duas iterações. O limiar 255 encerra a faixa; a sequência de passo 16 termina em 240. Serão mantidos kernel elíptico de $3\times3$, ausência de suavização adicional e filtro de área de 3 a 300 pixels. Uma amostra de 48 quadros anotados por vídeo de treino será selecionada por posições igualmente espaçadas no conjunto ordenado dos índices elegíveis; 12 desses quadros por vídeo formarão o subconjunto comum da triagem, e um quadro desse subconjunto será reservado à calibração de custo. Os pixels serão obtidos por decodificação sequencial dos MP4 canônicos, sem redimensionamento ou nova compressão, com hashes das entradas e das anotações.

Todas as configurações serão comparadas nos mesmos 144 quadros, usando a média dos F1 agregados dentro de cada um dos 12 vídeos. Desempates considerarão revocação, menor MAE de contagem e identificador, nessa ordem; tempo de execução não participa da escolha. Os cinco primeiros candidatos de uma bateria completa seguirão para um refinamento com limiares inteiros entre menos 15 e mais 15 em torno de cada candidato, mantendo sua morfologia e usando os 48 quadros por vídeo. Configurações repetidas serão deduplicadas, com limite de 155 candidatos. As duas finalistas desse refinamento poderão ser comparadas nos vídeos completos de validação. Essas etapas constituem seleção no treino e na validação; seus resultados não estimam desempenho em dados independentes. O protocolo registra orçamento e critérios antes da execução e interrompe comparações incompletas, sem classificar somente os candidatos que terminaram.

Antes de executar o refinamento, será realizado um novo benchmark com todos os candidatos nos mesmos 12 quadros reservados à medição de custo. A projeção somará o tempo de validação das entradas a 96 vezes o tempo do laço desse benchmark: fator 48 pela ampliação da amostra e fator 2 como margem operacional. O limite de liberação será de 4.800 segundos, com teto de 7.200 segundos para a bateria de refinamento. A projeção não constitui garantia de pior caso temporal. Todos os candidatos serão reavaliados nos 576 quadros; não serão somadas métricas antigas a resultados parciais novos. A escolha das duas finalistas seguirá somente a ordenação registrada, sem exigir diversidade de morfologia. Como a amostra passa de 12 para 48 quadros por vídeo, diferenças entre os valores agregados dessas duas etapas não serão interpretadas isoladamente como melhora do detector.

Em 8 de setembro de 2026, a triagem e o refinamento do threshold v3 estavam concluídos e conferidos. O refinamento avaliou 117 configurações nos mesmos 576 quadros de treino, resultando em T219/o0/c2 e T218/o0/c2 como finalistas, sem promoção do detector ou conclusão sobre a hipótese de fluxo. O protocolo da bateria subsequente previa comparar somente essas duas configurações nos quatro vídeos completos de validação: 14, 19 e 36, com 1.470 quadros cada, e 52, com 1.440, totalizando 5.850 quadros físicos e 11.700 avaliações de configuração--quadro. A seleção prevista seguia F1 macro por vídeo a 10 pixels, revocação, menor MAE de contagem e identificador; sensibilidades de 15/20 pixels e custo tinham caráter descritivo. Não estavam previstos testes de hipótese ou intervalos de confiança nessa etapa de seleção, nem congelamento automático ou acesso ao teste. A conclusão posterior dessa bateria é registrada abaixo.

O contrato anterior a essa execução exigia código e protocolo registrados em commit, com controles testados para exigir a quantidade exata de quadros e detectar quadros extras, rejeitar anotação ausente ou malformada e conferir hashes das entradas e das configurações. Esses controles ainda estavam em correção na revisão anterior à bateria. O orçamento prospectivo então fixado era de 1.800 segundos de tempo de parede, com controle operacional suave, limites de 2 GiB de RSS e de artefatos, e 2.000 previsões por quadro. A ordem dos oito pares configuração--vídeo seria embaralhada com seed 42, sem interpretar a ordem como repetição independente, e não seriam gerados vídeos de saída. Falhas deveriam preservar os artefatos parciais e impedir a seleção até a conclusão e conferência de todos os pares.

As detecções serão associadas temporalmente por métodos de complexidade crescente: centroide guloso, associação Húngara global, SORT, que combina Filtro de Kalman e associação Húngara \cite{bewley2016sort}, e uma implementação ByteTrack-style que reaproveita detecções de baixa confiança \cite{zhang2022bytetrack}. Uma variante Adaptive Flow-SORT será tratada separadamente como híbrida. Cada tracker será avaliado primeiro com caixas do ground truth, isolando a associação, e depois com o detector congelado, medindo o cenário fim a fim. Limiares de associação e o número de quadros tolerados sem detecção serão escolhidos na validação. As saídas terão a forma $(v,t,i,x,y)$, em que $v$ identifica o vídeo, $t$ o quadro, $i$ a célula e $(x,y)$ a posição de seu centro.

Antes da avaliação de identidades e da construção de janelas, será especificada a elegibilidade das trajetórias individuais das classes 0 e 2, com tratamento explícito das anotações de agrupamento, mudanças de classe e lacunas. Uma trajetória de caixa de agrupamento não será apresentada como trajetória de uma única célula. O protocolo de associação temporal e a preparação das entradas para o HOTA oficial serão documentados separadamente, sem transferir automaticamente a tolerância de 10 pixels usada na detecção.

Em seguida, a validação completa foi executada sob o commit \texttt{7f47afb}, com Git limpo e 824 testes aprovados antes da bateria. As oito runs somaram 11.700 avaliações em 166,657535 segundos, com pico amostrado de RSS de 236,594 MiB. T218/o0/c2 foi selecionada pelo critério prospectivo, com F1 macro de 0,658959, revocação macro de 0,794711 e MAE de contagem macro de 4,408433. T219/o0/c2 obteve F1 macro de 0,657819. A diferença de 0,11403 ponto percentual é descritiva: T218 teve F1 maior em 19/36, e T219 em 14/52. A conferência independente reconstruiu os 11.700 registros e agregados, com 1.772.878 comparações de campos/valores e 144 verificações amostrais de matching com SciPy. Duas tentativas iniciais foram preservadas após correções de suposições do verificador sobre formato de metadados; nenhuma run científica foi alterada. Na conclusão da bateria, a seleção ainda não estava congelada e não constituía teste independente, avaliação de folds ou evidência sobre a contribuição do fluxo. Posteriormente, T218/o0/c2 foi fixada como linha de base de desenvolvimento, com parâmetros e avaliação preservados. Esse escopo mantém teste, folds e aplicação bloqueados até um desenho confirmatório próprio, sem restaurar a cegueira historicamente perdida.

\subsection{Busca entre famílias de detectores: 11 de setembro de 2026}
\label{sec:busca-classicos-20260911}

A comparação entre famílias foi iniciada por uma busca registrada no treino, sob o commit \texttt{2547109}. Foram avaliadas 43 configurações em seis famílias: threshold fixo T218/o0/c2 como referência de desenvolvimento, Otsu, threshold adaptativo gaussiano, realce com CLAHE seguido de threshold, Blob e Watershed. Todas receberam os mesmos 576 quadros anotados, 48 por vídeo nos 12 treinos, reutilizados do cache autenticado, totalizando 24.768 avaliações de configuração--quadro. A métrica principal foi a média dos F1 por vídeo a 10 pixels, após agregação das contagens dentro de cada vídeo, com sensibilidades de 15/20 pixels e avaliação secundária de todos os objetos. Quadros, configurações e observações repetidas não foram tratados como unidades independentes.

O primeiro smoke foi interrompido pelo teto de 2.000 previsões por quadro, sem seleção. O plano e a execução parcial foram preservados. Uma revisão operacional explícita elevou somente esse teto para 307.200 previsões por quadro, mantendo a grade científica, os dados, as métricas e os limites de 2 GiB de memória e de artefatos. Um novo smoke completo, com 516 avaliações, precedeu a busca. Essa revisão não alterou as regras da execução histórica nem truncou previsões para concluir a comparação.

Na grade avaliada, \texttt{blob\_v1\_003} obteve F1 macro de 0,791005 e T218/o0/c2 obteve 0,775634, diferença de aproximadamente 1,54 ponto percentual. Blob apresentou F1 maior em sete dos 12 vídeos, e T218 nos outros cinco. Portanto, a diferença média não representa superioridade uniforme. As famílias tiveram quantidades diferentes de configurações pesquisadas, e T218 já havia sido ajustado em uma busca anterior. O resultado descreve esta triagem limitada no treino, não o melhor desempenho possível de cada família ou sua generalização. Foram retidas duas finalistas por família pesquisada e a referência T218, totalizando 11 candidatas à próxima etapa, sem promoção, liberação de validação ou acesso ao teste.

A busca terminou em 752,160737 segundos, com RSS máximo amostrado de 823,145 MiB. A conferência independente verificou 289 arquivos, 30.927.965 comparações e 148.608 associações com SciPy, sem diferença numérica nos valores conferidos. A paridade de T218 nos mesmos quadros foi confirmada contra o refinamento anterior. Essa conferência autentica e reconstrói os derivados; não reexecuta os detectores nem a decodificação original. A comparação ainda não inclui YOLO, MOG2 ou KNN: o primeiro exige treinamento próprio, e os métodos temporais precisam de processamento contínuo e aquecimento, incompatíveis com tratar os quadros espaçados como imagens independentes. O refinamento e a validação das alternativas exigirão protocolos novos; maior F1 de detecção, por si só, não assegura maior HOTA ou menor erro de predição.


\subsection{Refinamento local e preparação dos dados aprendidos}
\label{sec:refinamento-classicos-20260911}

Após a busca anterior, foram registrados o contrato e a implementação de um refinamento local, antes das novas inferências. As duas finalistas de cada uma das cinco famílias pesquisadas originaram 54 propostas que incluíam os próprios pais e alterações de um parâmetro por vez. Três propostas ficaram fora dos limites prospectivos; as 51 válidas resultaram em 45 configurações únicas após deduplicação, preservando todas as origens. A referência T218 foi autenticada e não reexecutada nesta etapa.

Um smoke de 540 avaliações precedeu o refinamento completo. Ambos foram concluídos e conferidos independentemente sob o commit \texttt{ca68f16}. O refinamento utilizou os mesmos 576 quadros físicos dos 12 treinos, totalizando 25.920 avaliações. A melhor configuração da grade de cada família apresentou os seguintes F1 macro a 10 pixels: Blob: 0,791005; Watershed: 0,627005; Otsu: 0,626315; Adaptativo: 0,446885; Híbrido CLAHE: 0,364491. Essas médias descrevem seleção nos dados de treino, com esforço desigual entre famílias. Foram retidas duas finalistas por família, totalizando dez, sem promoção, abertura do teste ou escolha final da pipeline.

Os dez pais reproduziram suas detecções, anotações exportadas e métricas científicas nos mesmos quadros, exceto tempo e identificação da nova execução. A conferência independente reconstruiu vizinhança, associações, agregações e classificação. Esse controle verifica consistência e reprodutibilidade dos derivados; não repete a inferência nem certifica generalização. A duração da bateria foi 528,244100 segundos e o máximo de RSS amostrado foi 410,406000 MiB. Recursos e custos permanecem vinculados ao ambiente e ao commit da execução.

Foi também preparado um dataset derivado independente para o futuro treinamento YOLO: 17.466 pares de imagem e anotação no treino e 5.850 na validação, totalizando 23.316. As 174 lacunas do vídeo 23 foram excluídas, sem tratá-las como negativos. As três classes foram preservadas; os rótulos geométricos YOLO foram confrontados com as anotações FTID usadas na avaliação, ignorando somente o identificador. Fontes, cópias e referências foram autenticadas por hashes, e as cópias JPEG foram decodificadas e conferidas independentemente. Não houve leitura de imagens do teste ou treinamento de rede nesta preparação.

O dataset preparado permanece imutável. O futuro executor de treinamento deverá usar outra cópia independente para isolar os caches e eventuais reparos da biblioteca, mantendo rastreável qualquer mudança de entrada. Receita, pesos iniciais, sementes, seleção de checkpoints e parâmetros de inferência exigem contrato próprio. A sequência seguinte é validar as finalistas clássicas em vídeos completos e treinar e comparar YOLO sob a mesma regra de detecção; os métodos temporais requerem protocolo contínuo com aquecimento.

A paridade geométrica dos rótulos não certifica igualdade de pixels entre os JPEGs recebidos e os quadros decodificados dos vídeos. Os JPEGs constituem as entradas de treinamento preparadas; a comparação de detecção entre YOLO e os clássicos deverá processar os mesmos quadros do vídeo ou cache de avaliação, com pré-processamento explicitamente registrado.

\section{Estimação do movimento aparente}

O movimento aparente será estimado por Lucas--Kanade, Farnebäck \cite{farneback2003twoframe}, Horn--Schunck e, se os recursos computacionais permitirem, RAFT \cite{teed2020raft}. Variantes híbridas combinarão compensação de movimento global, consistência direto--reverso, rejeição robusta e eventual refinamento neural. Regiões ocupadas pelas células e por artefatos evidentes serão mascaradas quando o objetivo for caracterizar o fundo; o clássico puro e cada adaptação manterão configuração e resultado próprios.

Como os dados do VISEM e as anotações do VISEM-Tracking não fornecem uma medição física do campo de velocidade do fluido, o resultado será tratado como campo de movimento aparente. Sua consistência será analisada por erro fotométrico após a deformação de um quadro em direção ao seguinte, concordância entre fluxos direto e reverso, regularidade temporal e estabilidade nas regiões de fundo. Sequências sintéticas com deslocamentos conhecidos poderão complementar a verificação controlada dos estimadores, sem constituir uma segunda coleção de vídeos microscópicos reais. A conversão dos deslocamentos aparentes para unidades físicas somente será realizada quando houver calibração espacial e temporal confiável; essa conversão, por si só, não estabelece a velocidade física do fluido.


Como verificação inicial da conexão causal com a referência individual, foi
registrado e executado um ensaio de Farnebäck sob o commit \texttt{16eecbb},
após 1.495 testes automatizados. Foram utilizados somente os quadros de 0 a 19
dos vídeos de treino 11 e 12, sem máscara ou busca de parâmetros. Na origem
$t=19$, cada um dos 19 campos históricos $F_s$, que representa $s\to s+1$,
foi amostrado no centro anotado $\mathbf p_s$ do primeiro quadro do par, com
interpolação bilinear e pesos em precisão de 64 bits. O último par disponível
é $18\to19$; nenhuma imagem posterior foi fornecida ao estimador. O consumidor
recebeu apenas o histórico e as chaves das janelas, sem coordenadas futuras.
A elegibilidade continua condicionada à disponibilidade de futuro completo
na referência derivada, exclusivamente para avaliação retrospectiva.

As 68 janelas selecionadas antes dos pixels produziram 1.292 amostras com
suporte numérico válido, sem exclusões. Os 38 pares foram estimados nos dois
sentidos, totalizando 76 campos, além de 36 comparações temporais. Os erros
fotométricos médios após alinhamento foram 0,002374 e 0,003461 nos vídeos 11
e 12, contra 0,003008 e 0,005554 com deslocamento zero, em intensidade cinza
dividida por 255 e sobre o mesmo suporte de cada par. As consistências
direto--reverso médias foram 0,015067 e 0,169551 pixels. As médias por vídeo
dão peso igual aos pares com suporte. Validade numérica de todas as amostras
não implica acurácia de 100\%, e esses diagnósticos não medem a velocidade
física do fluido nem o ganho em predição.

O ensaio levou 56,923953 segundos, com RSS máximo amostrado de 289,598 MiB.
A conferência independente verificou 190 arquivos e 37.406 comparações,
das quais 8.378 numéricas, com diferença máxima de $8{,}88\times10^{-16}$,
abaixo da tolerância absoluta de $10^{-9}$. Uma falha inicial na leitura do
esquema do relatório histórico da referência foi preservada e corrigida sob
\texttt{0fccda8}, após nove testes sintéticos, sem modificar ou repetir a run.
A conferência reconstituiu a matemática dos derivados, sem importar a
implementação científica; não constitui decodificação independente do MP4 ou
reexecução do estimador. O ensaio certifica a conexão operacional nesse
prefixo de dois vídeos. A extração ampliada e a ablação pareada com e sem fluxo
permanecem sujeitas a contrato prospectivo próprio.


Para preparar a escala, foi registrado antes dos pixels um benchmark de
armazenamento compacto sob o commit \texttt{6110c44}, após 1.728 testes.
Foram processados os quadros de 0 a 59 de cada um dos 12 vídeos de treino,
com a mesma configuração de Farnebäck, uma thread de CPU e OpenCL desativado.
Os 708 campos diretos foram amostrados nas posições históricas de origem.
Foram retidos campos densos nos dois sentidos apenas nos pares $0\to1$ e
$58\to59$ de cada vídeo, além de seus quadros cinza. As 10.848 janelas
compartilham 15.762 amostras distintas, correspondentes a 206.112 usos.
Cada amostra preserva os quatro vizinhos usados na interpolação; o QA pode
reconstruir seus pesos e valores, mas somente nos sentinelas pode confrontar
os vizinhos com o campo denso retido. Não houve segunda decodificação ou
reexecução independente do estimador. Todas as amostras tiveram suporte
numérico válido, sem que isso certifique acurácia do movimento estimado.

O benchmark terminou em 199,583442 segundos, com RSS amostrado de 375,145 MiB
e 154.023.508 bytes finais. A conferência independente passou na primeira
tentativa: 288 arquivos e 934.668 comparações, com diferença máxima de
$3{,}55\times10^{-15}$. A regra de projeção, fixada antes dos resultados e
com fator de segurança dois, estimou 8.052,675671 segundos e 1.768,62 MiB para
a escala completa. O tempo excedeu o teto prospectivo de 7.200 segundos;
o armazenamento permaneceu abaixo dos 4.096 MiB. Por isso, a extração completa
não foi liberada por aquele protocolo. Na conclusão do benchmark, a continuidade
prevista era discriminar custos de decodificação, estimação e amostragem antes
de ampliar a extração. Em 11 de setembro, o tempo deixou de ser impedimento
para novos planos: o teto e a projeção históricos permanecem preservados, mas
uma nova execução pode monitorar tempo sem aquele corte, mantendo controle de
memória, armazenamento e completude. A instrumentação deixa de ser condição
obrigatória para a ablação exploratória nas 10.848 janelas já disponíveis,
que ainda exige contrato e executor próprios. Não houve avaliação de ADE/FDE
com fluxo, seleção de parâmetros ou teste da hipótese neste benchmark.

\section{Predição das trajetórias}

Antes dos preditores, foi registrado e executado o protocolo de referência individual apenas nos 12 vídeos de treino, sob o commit \texttt{33d191d}, com Git limpo e 1.026 testes aprovados. As classes 0 e 2 preservam o ID original; classe 1, ausência do ID e ausência de anotação interrompem os segmentos. Mudanças entre classes 0 e 2 não interrompem o trecho. A segmentação não renumera identidades biológicas e não define implicitamente a política de HOTA. Todas as observações e todos os segmentos são mantidos, inclusive os curtos.

Foram preservadas 363.074 observações individuais em 725 segmentos, dos quais 623 geram 343.776 janelas de 20 posições observadas e dez posições futuras, com passo de um quadro. Os 102 segmentos sem janela conservam 1.231 observações. Das origens individuais, 12.914 não têm histórico completo e 6.384 têm histórico, mas não o futuro completo. A elegibilidade é condicionada à disponibilidade de referência individual futura; esse uso é exclusivo dos alvos e da avaliação retrospectiva, não das entradas do preditor. As janelas permanecem dependentes dentro do vídeo. A conferência independente reconstruiu todas as observações, segmentos e janelas a partir dos derivados, verificando 86 arquivos e 9.407.090 comparações. A preparação levou 151,929874 segundos, com pico amostrado de RSS de 226,148 MiB, sem decodificar pixels, avaliar modelos ou abrir fontes de validação e teste. Não houve transição direta de um ID individual para classe 1 nesta referência de treino; a regra foi verificada sinteticamente e não permite inferir a causa biológica de ausências. Em etapa posterior, os índices foram consumidos pelos baselines fixos descritos a seguir, com entradas causais e avaliação densa de ADE.


Como primeira avaliação de predição, persistência e velocidade constante foram fixadas antes dos resultados, sob o commit \texttt{5289c93}. A persistência repete a última posição; a velocidade constante utiliza a mediana componente a componente das cinco últimas diferenças, calculadas a partir de seis posições. Ambos receberam apenas o histórico, em precisão de 64 bits, e produziram todos os passos de 1 a 10 sem limitar as previsões aos contornos da imagem. As 343.776 janelas foram avaliadas por cada método, cobrindo 606 dos 669 IDs individuais originais. Os 63 IDs sem janela permanecem na descrição de cobertura, sem imputação de erro zero. A agregação principal calcula a média de janelas por ID original, reunindo seus segmentos, seguida da média com peso igual entre IDs dentro do vídeo e entre os 12 vídeos. A média por janela dentro de vídeo, seguida de pesos iguais entre vídeos, é secundária.

Na agregação principal, a persistência obteve ADE de 3,890964 pixels e FDE de 6,619751 pixels no horizonte de dez quadros; a velocidade constante obteve 3,137020 e 5,774373 pixels, respectivamente. A redução não foi uniforme: a velocidade constante apresentou ADE menor em nove vídeos e FDE menor em sete. São resultados descritivos nas trajetórias anotadas de treino, condicionados à disponibilidade de futuro de referência, sem ajuste de hiperparâmetros, seleção de método, inferência estatística ou confirmação da hipótese de fluxo. O FPS nominal foi 48 no vídeo 35, 50 no vídeo 82 e 49 nos demais; portanto, dez quadros não correspondem exatamente à mesma duração física em todos os vídeos. Nenhuma reamostragem foi aplicada.

Antes da bateria, 1.362 testes automatizados foram aprovados. A execução completa levou 54,191753 segundos, com RSS amostrado de 199,027 MiB. Uma conferência independente posterior reconstituiu todas as previsões, os erros densos e as agregações a partir dos arquivos exportados e da referência derivada, sem importar a implementação científica ou abrir fontes originais. O manifesto, os arquivos por janela e por ID, a comparação por vídeo e a figura preservam a proveniência da execução. Fluxo, rastreamento e modelos aprendidos não foram avaliados nessa bateria.

Cada trajetória válida será dividida em histórico de 20 quadros e horizontes de 1, 5 e 10 quadros, sem atravessar vídeo, identidade, split ou lacuna de anotação. As entradas básicas serão posição, deslocamento e velocidade estimada. As configurações com fluxo acrescentarão o vetor aparente amostrado na posição da célula ou em um anel válido ao seu redor. Serão comparados persistência, velocidade constante, Filtro de Kalman, Filtro de Partículas e LSTM, além de variantes híbridas correspondentes com fluxo. A LSTM foi escolhida por representar dependências temporais e já ter sido aplicada à previsão da localização de espermatozoides \cite{noy2023location}.

O experimento principal será uma ablação pareada. O mesmo preditor, com os mesmos dados, janelas, hiperparâmetros e sementes, será treinado primeiro sem fluxo e depois com características de fluxo. A variável controlada será, portanto, a presença do movimento aparente, e não uma mudança simultânea de arquitetura. A avaliação ocorrerá em dois cenários: trajetórias anotadas, para isolar o preditor, e trajetórias produzidas pelo rastreador, para medir a degradação da pipeline completa.

Essa ablação exigirá um manifesto comum de janelas, identificadas por vídeo, identidade, instante final do histórico e horizonte. Os dois métodos receberão as mesmas janelas elegíveis; serão informadas a cobertura, as exclusões e as falhas por método. Ausência de fluxo válido não será substituída por vetor zero, nem poderá excluir silenciosamente somente os casos difíceis de um preditor. Executar métodos separadamente não certifica esse pareamento, que será verificado antes de produzir a comparação científica.

Para uma previsão realizada no instante $t$, todas as imagens, máscaras e características de entrada deverão usar somente informações disponíveis até $t$. O fluxo entre $t-1$ e $t$ é admissível; fluxo que utilize imagens posteriores a $t$ não poderá entrar na avaliação principal. O futuro será reservado aos alvos. Qualquer controle que forneça informação futura será identificado como oráculo e relatado separadamente. Normalizações e parâmetros estimados a partir dos dados serão ajustados somente no treino correspondente.


Foi também implementada uma variante causal fixa da velocidade constante,
em precisão de 64 bits. Para as cinco últimas transições observadas, define-se
$\mathbf g_s=F_s(\mathbf p_s)$ e utiliza-se
\begin{equation}
\widehat{\mathbf p}_{t+h}=\mathbf p_t+h\left[
\mathop{\mathrm{mediana}}_{s=t-5}^{t-1}
(\mathbf p_{s+1}-\mathbf p_s-\mathbf g_s)+\mathbf g_{t-1}\right],
\quad h=1,\ldots,10.
\end{equation}
A mediana é calculada por componente. O resíduo é algébrico, sem interpretação
de velocidade intrínseca ou física do fluido. Se o fluxo histórico é constante,
sua subtração e reposição reproduzem o baseline de velocidade constante dentro
da tolerância numérica; esse controle e as recusas de entradas futuras foram
verificados sinteticamente. A avaliação real permanece pendente. Os dois braços
deverão usar a mesma coorte com as cinco amostras finais válidas, reportando
também a cobertura das 19 transições. Ausências anteriores que não entram na
fórmula não excluem por si mesmas a janela, e ausência de fluxo não é imputada
como vetor zero.

\section{Critérios de validação e comparação}

Na detecção, a métrica de promoção será o F1 espacial descrito anteriormente. No rastreamento, a métrica principal será HOTA, acompanhada por MOTA, IDF1, trocas de identidade e fragmentações \cite{luiten2021hota}. Na predição serão calculados ADE e FDE \cite{alahi2016sociallstm}. Os erros serão expressos em pixels e, quando houver calibração confiável, também em micrômetros.

Para cada horizonte $H$, a definição operacional utilizará todos os passos futuros de 1 a $H$:
\begin{equation}
\mathrm{ADE}_{H}=\frac{1}{H}\sum_{h=1}^{H}\left\|\widehat{\mathbf p}_{t+h}-\mathbf p_{t+h}\right\|_2,
\qquad
\mathrm{FDE}_{H}=\left\|\widehat{\mathbf p}_{t+H}-\mathbf p_{t+H}\right\|_2.
\end{equation}
Os horizontes 1, 5 e 10 serão pontos de apresentação; cada ADE exigirá os erros intermediários correspondentes. A média apenas dos erros nos passos 1, 5 e 10 não será denominada ADE de 1 a 10. A interface histórica de avaliação calcula a média dos instantes fornecidos e permanece identificada como tal. A bateria fixa de persistência e velocidade constante utilizou uma interface densa própria, verificada com casos sintéticos e por reconstrução independente dos resultados. As futuras variantes deverão satisfazer o mesmo contrato denso e a separação entre histórico e alvos.

As comparações abrangerão linhas de base clássicas, métodos aprendidos e híbridos, preditores sem e com fluxo e trajetórias anotadas e estimadas. Para evitar que vídeos longos dominem a média, frames e trajetórias serão agregados primeiro dentro de cada vídeo, que será a unidade independente. Friedman para três ou mais métodos e Wilcoxon pareado, com Holm para a família de comparações planejada, somente serão interpretados conforme a adequação do desenho e do tamanho amostral. Serão apresentados os valores e as diferenças pareadas por vídeo, além de média e mediana dos efeitos. Treinamentos estocásticos serão repetidos com sementes distintas, agregadas dentro do vídeo. Qualidade e custo serão apresentados por fronteiras de Pareto, sem um placar único com pesos subjetivos.

No holdout de quatro vídeos, o menor p possível para o Wilcoxon bilateral exato, com diferenças não nulas e sem empates em seus módulos, é $2/2^4=0{,}125$. Assim, esse teste não permite prometer significância a 5\% nesse conjunto. Intervalos de 95\% por bootstrap agrupado de quatro vídeos, caso apresentados, terão caráter exploratório e acompanharão os valores individuais; frames, trajetórias e seeds não ampliarão artificialmente o número de unidades independentes. Uma análise fora da amostra dos 20 vídeos dependerá de seleção e ajuste dentro dos folds de treino de cada avaliação externa, para limitar o viés de seleção\footnote{Cawley e Talbot (2010): \url{https://www.jmlr.org/papers/volume11/cawley10a/cawley10a.pdf}.}. Esse desenho não elimina a necessidade de declarar a exposição histórica do teste.

Os resultados de detecção e rastreamento no VISEM-Tracking serão confrontados com os protocolos publicados por Thambawita et al. \cite{thambawita2023visemtracking} e Zhang et al. \cite{zhang2024spermyolo}. A predição será comparada à linha de base de velocidade constante e ao uso de LSTM descrito por Noy et al. \cite{noy2023location}. Diferenças de divisão dos dados, resolução, horizonte, hardware e implementação serão explicitadas para evitar comparações diretas entre condições incompatíveis.

\section{Etapas de desenvolvimento}

A estratégia inicial priorizou uma pipeline mínima para isolar a hipótese, com referência individual, baselines de persistência e velocidade constante e conexão causal do fluxo. Esses baselines e os derivados de fluxo já foram produzidos e conferidos nas etapas descritas; a ablação real ainda não ocorreu. Na continuidade de 11 de setembro, a comparação e a validação das alternativas de detecção e rastreamento passaram a compor uma frente experimental explícita, iniciada pela busca da Seção~\ref{sec:busca-classicos-20260911}. O refinamento clássico e a preparação autenticada do dataset YOLO foram concluídos, conforme a Seção~\ref{sec:refinamento-classicos-20260911}. As finalistas clássicas ainda exigem validação completa; YOLO precisa de treinamento e comparação no protocolo atual. O contrato de rastreamento deve ser concluído antes da avaliação por HOTA, distinguindo associação com caixas anotadas sem IDs e associação com detecções automáticas.

A escolha da combinação para rastreamento e predição deverá decorrer dessas comparações e da avaliação da cadeia completa, preservando seus erros de localização, identidade e cobertura. Em paralelo, a ablação exploratória nas janelas GT com fluxo já extraído poderá isolar a contribuição dessa característica, sem ser apresentada como resultado do sistema automático. Os modelos aprendidos e seus controles de capacidade deverão ser avaliados antes da síntese final. Um resultado sem redução de ADE/FDE também será informativo, desde que a comparação seja causal, pareada e reproduzível. O roteiro a seguir representa a sequência global do projeto; as etapas já concluídas não serão reiniciadas.

\begin{enumerate}
	\item \textbf{Auditoria e padronização dos dados.} Os vídeos, metadados e anotações serão inventariados, convertidos para formatos comuns e verificados visualmente.

	\item \textbf{Definição das divisões e linhas de base.} Serão fixados os conjuntos de treinamento, validação e teste, as sementes e as configurações clássicas de referência.

	\item \textbf{Treinamento e avaliação da detecção.} A abordagem clássica e o detector aprendido serão ajustados e comparados sob os mesmos dados.

	\item \textbf{Implementação do rastreamento.} Os métodos por centroide, SORT e ByteTrack serão aplicados às mesmas detecções para separar falhas de localização e associação.

	\item \textbf{Estimação e análise do fluxo.} Farnebäck e, se viável, RAFT serão executados com compensação de movimento global e máscaras de primeiro plano.

	\item \textbf{Construção dos conjuntos de predição.} As trajetórias serão segmentadas por continuidade e elegibilidade documentadas, preservando as exclusões. Janelas observadas e futuras serão separadas; normalizações dos preditores serão ajustadas somente no treino correspondente, sem filtrar trajetórias por velocidade ou erro futuro.

	\item \textbf{Treinamento e ablação dos preditores.} As configurações com e sem fluxo serão avaliadas nos cenários controlados.

	\item \textbf{Análise e documentação.} Os resultados serão agregados, acompanhados por intervalos de confiança, visualizações e exemplos de falha, e relacionados à hipótese de pesquisa.
\end{enumerate}

\section{Declaração sobre o uso de Inteligência Artificial}

Para a consolidação desta versão inicial, utilizei as ferramentas ChatGPT e Codex, da OpenAI, como apoio à organização dos conteúdos das entregas anteriores, à pesquisa bibliográfica, à revisão linguística, à conferência da coerência entre os capítulos e à formatação em LaTeX. O tema, os objetivos, as decisões metodológicas e a revisão final do texto são de minha responsabilidade, e a submissão será realizada após a concordância do orientador.
